% ============ 22 Künneth 公式 (原书 pp.135-137) ============
\chapter{K\"unneth 公式}\label{sec:22}
\index{K\"unneth 公式}
\index{杯积}
\index{拟同构}
\index{投影公式}
\index{Tor 群}

这只是简要总结.形式上，K\"unneth 公式的理论对平展拓扑与拓扑空间相同---只需把容易的拓扑紧合基变换定理换成困难得多的平展版本.拓扑版本见 Iversen 1986, VII.2.

\textbf{杯积} 设 $F$ 与 $G$ 是 $X$ 上的层.则有杯积映射
\[
H^{r}(X_{\mathrm{et}}, F) \times H^{s}(X_{\mathrm{et}}, G) \to H^{r+s}(X_{\mathrm{et}}, F \otimes G).
\]
定义它们最简单的方法是把这些群等同于 \v{C}ech 上同调群，并令
\[
(f \smile g)_{i_0 \cdots i_{r+s}} = f_{i_0 \cdots i_r} \otimes g_{i_r \cdots i_{r+s}}.
\]

\subsection*{K\"unneth 公式}

设 $\Lambda$ 是有限环.$\Lambda$-模复形 $A^{\bullet} \to B^{\bullet}$ 的映射称为拟同构，如果它在复形的上同调群上诱导同构.

设 $X$ 与 $Y$ 是代数闭域 $k$ 上的代数簇，$F$ 与 $G$ 分别是 $X$ 与 $Y$ 上的层.考虑
\[
\begin{tikzcd}
X \times Y \arrow[rd, "p"'] \arrow[dd, "q"] & \\
& Y \\
X &
\end{tikzcd}
\]
把限制映射
\[
H^{r}(X, F) \to H^{r}(X \times Y, p^{*}F)
\]
\[
H^{s}(Y, G) \to H^{s}(X \times Y, q^{*}G)
\]
与杯积映射
\[
H^{r}(X \times Y, p^{*}F) \times H^{s}(X \times Y, q^{*}F) \to H^{r+s}(X \times Y, p^{*}F \otimes q^{*}G)
\]
结合，得到映射
\[
H^{r}(X, F) \times H^{s}(Y, G) \to H^{r+s}(X \times Y, p^{*}F \otimes q^{*}G).
\]
K\"unneth 公式研究映射
\[
\bigoplus_{r+s=n} H^{r}(X, F) \otimes H^{s}(Y, G) \to H^{n}(X \times Y, p^{*}F \otimes q^{*}G)
\]
在多大程度上是同构.

\begin{theorem}\label{thm:22.1}
设 $X$ 与 $Y$ 是完备簇；存在拟同构
\[
\mathcal{H}(X, \Lambda) \otimes_{\Lambda} \mathcal{H}(Y, \Lambda) \to \mathcal{H}(X \times Y, \Lambda),
\]
其中 $\mathcal{H}(X, \Lambda)$，$\mathcal{H}(Y, \Lambda)$ 与 $\mathcal{H}(X \times Y, \Lambda)$ 是 $\Lambda$-模复形，使得
\begin{enumerate}
  \item $H^{r}(\mathcal{H}(X, \Lambda)) \cong H^{r}(X_{\mathrm{et}}, \Lambda)$，另外两个类似；
  \item $\mathcal{H}(X, \Lambda)$ 是平坦 $\Lambda$-模复形.
\end{enumerate}
\end{theorem}

\begin{corollary}\label{cor:22.2}
有谱序列
\[
\bigoplus_{i+j=s} \operatorname{Tor}^{\Lambda}_{r}(H^{i}(X_{\mathrm{et}}, \Lambda), H^{j}(Y_{\mathrm{et}}, \Lambda)) \Rightarrow H^{r+s}(X \times Y, \Lambda).
\]
\end{corollary}

\begin{proof}
定理中的拟同构产生这样的谱序列是同调代数的标准结果.
\end{proof}

\begin{remark}\label{rem:22.3}
(a) 遗憾的是，一般地，推论中的谱序列是无限的.例如，设 $\Lambda = \mathbb{Z}/\ell^{2}\mathbb{Z}$.有任意长度的正合序列
\[
0 \to \mathbb{Z}/\ell\mathbb{Z} \to \mathbb{Z}/\ell^{2}\mathbb{Z} \to \cdots \to \mathbb{Z}/\ell^{2}\mathbb{Z} \xrightarrow{\ell} \mathbb{Z}/\ell^{2}\mathbb{Z} \to \mathbb{Z}/\ell\mathbb{Z} \to 0.
\]
这表明 $\operatorname{Tor}^{\mathbb{Z}/\ell^{2}\mathbb{Z}}_{r}(\mathbb{Z}/\ell\mathbb{Z}, \mathbb{Z}/\ell\mathbb{Z})$ 对无限多个 $r$ 非零.

(b) 定理的证明(见下面)使用紧合基变换定理.它对非完备簇与紧支集上同调也成立，若簇光滑或假定 Deligne 的定理(下面的 24.3)，则对通常上同调也成立.

(c) 要陈述和证明定理，最自然的是使用导出范畴.

(d) 定理对 $\Lambda$-模的可构造层 $F$ 与 $G$ 成立，只要其中至少一个是平坦 $\Lambda$-模层(即 $F \otimes_{\Lambda} \cdot$ 正合).

对 $\Lambda = \mathbb{Z}/\ell^{n}\mathbb{Z}$，可以相容地定义复形与拟同构，使 $n$ 变动，从而可以取逆向极限，得到拟同构
\[
\mathcal{H}(X, \mathbb{Z}_{\ell}) \otimes_{\mathbb{Z}_{\ell}} \mathcal{H}(Y, \mathbb{Z}_{\ell}) \to \mathcal{H}(X \times Y, \mathbb{Z}_{\ell}).
\]
因为 $\mathbb{Z}_{\ell}$ 是主理想整环，这给出通常形式的 K\"unneth 公式.
\end{remark}

\begin{theorem}\label{thm:22.4}
对簇 $X$ 与 $Y$，有典范正合序列
\[
0 \to \bigoplus_{r+s=m} H^{r}(X, \mathbb{Z}_{\ell}) \otimes H^{s}(Y, \mathbb{Z}_{\ell}) \to H^{m}(X \times Y, \mathbb{Z}_{\ell})
\]
\[
\to \bigoplus_{r+s=m+1} \operatorname{Tor}^{\mathbb{Z}_{\ell}}_1(H^{r}(X, \mathbb{Z}_{\ell}), H^{s}(Y, \mathbb{Z}_{\ell})) \to 0
\]
(对紧支集上同调类似).
\end{theorem}

\subsection*{证明}

第一步是证明投影公式.

\begin{proposition}\label{prop:22.5}[投影公式]
考虑正则映射 $f \colon X \to S$.则对 $X$ 上 $\Lambda$-模的任何平坦层 $F$ 与 $S$ 上有界上复形 $G^{\bullet}$，有拟同构
\[
(Rf_{*}F) \otimes G^{\bullet} \xrightarrow{\sim} Rf_{*}(F \otimes f^{*}G^{\bullet}).
\]
\end{proposition}

\begin{proof}
为证明它，把一般情形归结为 $G^{\bullet} = \Lambda$ 的情形，后者显然.
\end{proof}

其次，考虑
\[
\begin{tikzcd}
X \times_{S} Y \arrow[rd, "p"'] \arrow[dd, "q"] & \\
& Y \arrow[dd, "g"] \\
X \arrow[rd, "f"'] & \\
& S
\end{tikzcd}
\]
设 $F$ 是 $X$ 上 $\Lambda$-模的平坦可构造层，$G$ 是 $Y$ 上 $\Lambda$-模的可构造层.则在导出范畴的语言中，
\[
Rf_{*}F \otimes Rg_{*}G
\]
\[
\parallel \text{(投影公式)}
\]
\[
Rf_{*}(F \otimes f^{*}Rg_{*}G)
\]
\[
\parallel \text{(基变换)}
\]
\[
Rf_{*}(F \otimes Rp_{*}(q^{*}G))
\]
\[
\parallel \text{(投影公式)}
\]
\[
Rf_{*}(Rp_{*}(p^{*}F \otimes q^{*}G))
\]
\[
\parallel (Rf_{*} \circ Rp_{*} = R(f \circ p)_{*})
\]
\[
Rh_{*}(p^{*}F \otimes q^{*}G).
\]
为了能应用基变换定理，需要
\begin{enumerate}
  \item $g$ 紧合(用紧合基变换定理)，或
  \item $f$ 光滑(用光滑基变换定理)，或
  \item $S$ 是域的谱(用 Deligne 定理).
\end{enumerate}
